\ifdefined\iscombined
	\quiz{Quiz 5}
\else
\documentclass{article}
\usepackage{xparse}
\usepackage{tabularx}
\usepackage{changepage}
\usepackage{listings}
\usepackage{amsthm}
\usepackage{comment}
\usepackage{amsmath}
\usepackage{braket}
\usepackage{amsfonts}
\usepackage{amssymb}
\usepackage{sectsty}
\usepackage{hyperref}
\usepackage{fancyhdr}
\usepackage[top=1in,bottom=1in,left=1in,right=1in]{geometry}
\usepackage{float}
\usepackage{graphicx}
\usepackage{mathtools}
\usepackage{tikz}
\usepackage{enumitem}
\usetikzlibrary{arrows, arrows.meta}

\date{
	\vspace{-.75em}
	{\large \today}
}

\title{
	\vspace*{-1.5em}
	{\Huge Quiz 5} \\
	\vspace{-1em}
	\rule{\textwidth/2}{.01em} \\
	\vspace{-.75em}
}

\author{
	{\large Matthew Farah}
}

\hypersetup{
	colorlinks=true,
	linkcolor=blue,
	filecolor=magenta,
	urlcolor=blue,
	citecolor=blue,
	anchorcolor=blue
}

\fancyhead{}
\fancyhead[L]{\today}
\fancyhead[R]{Quiz 5}

\setlength{\parindent}{0cm}

\tikzset{vertex/.style = {shape = circle, draw, minimum size = 2em}}
\tikzset{edge/.style = {-,> = latex'}}

\newenvironment{solution}{
	\vspace{.5em}
	\par
	\begin{adjustwidth}{1.5em}{}
	\textbf{{Solution:}}
	\par
	\vspace{.5em}
}{
	\vspace{1em}
	\end{adjustwidth}
	\setcounter{step}{0}
}

\makeatother
\makeatletter

\newcounter{step}
\renewcommand{\thestep}{\Roman{step}}

\newenvironment{step}{
	\refstepcounter{step}
	\setcounter{substep}{0}
	\vspace{1em}\par\noindent
	\ignorespaces\noindent\textbf{(\thestep)}
	\ignorespaces
}{
	\par
	\vspace{1.5em}
}

\newcounter{substep}
\renewcommand{\thesubstep}{\alph{substep}}

\newenvironment{substep}{
	\refstepcounter{substep}
	\vspace{1em}\par\noindent
	\ignorespaces\noindent\textbf{(\thesubstep)}
	\ignorespaces
}{
	\par
	\vspace{1.5em}
}

\newcounter{problem}
\renewcommand{\theproblem}{\arabic{problem}}

\newenvironment{problem}[1][]{
	\newpage
	\refstepcounter{problem}
	\setcounter{subproblem}{0}
	\vspace{.5em}\par\noindent
	\begin{adjustwidth}{1.5em}{}
	\noindent\textbf{Problem \theproblem: }\noindent\textit{#1}
	\par
	\phantomsection
	\addcontentsline{toc}{section}{\protect{\large{Problem \theproblem}}}
}{
	\vspace{.5em}
	\end{adjustwidth}
}

\newcounter{subproblem}
\renewcommand{\thesubproblem}{\alph{subproblem}}

\newenvironment{subproblem}{
	\setcounter{step}{0}
	\refstepcounter{subproblem}
	\phantomsection
	\addcontentsline{toc}{subsection}{\protect{\large{(\thesubproblem)}}}
	\begin{adjustwidth}{1.5em}{}
	\noindent{\textbf{(\thesubproblem) }}\ignorespaces
}{
	\par
	\vspace{.5em}
	\end{adjustwidth}
}

\newcounter{lemma}
\renewcommand{\thelemma}{\arabic{lemma}}

\newenvironment{lemma}[1][]{
	\refstepcounter{lemma}
	\vspace{1em}\par\noindent
	\noindent\textbf{Lemma ({#1}): }\noindent\ignorespaces
}{
	\vspace{.5em}
}

\newcounter{proposition}
\renewcommand{\theproposition}{\arabic{proposition}}

\newenvironment{proposition}{
	\refstepcounter{proposition}
	\vspace{1em}\par\noindent
	\noindent\textbf{Proposition: }\noindent\ignorespaces
}{
	\vspace{.5em}
}

\newcommand{\supp}[1]{\text{supp}({#1})}
\newcommand{\ord}[1]{\text{ord}\left({#1}\right)}
\newcommand{\gen}[1]{\left\langle {#1} \right\rangle}
\newcommand{\lcm}[2]{\text{lcm}({#1}, {#2})}
\newcommand{\dprod}[2]{\mathbb{Z}_{#1} \times \mathbb{Z}_{#2}}

\renewcommand{\mod}[1]{\ (\text{mod} \; {#1})}

\sectionfont{\fontsize{12}{12}\bfseries}
\subsectionfont{\fontsize{10}{12}\normalfont}
\subsubsectionfont{\fontsize{10}{12}\normalfont}

\begin{document}

\maketitle
\pagestyle{fancy}

\newcolumntype{Y}{>{\centering\arraybackslash}X}

\tableofcontents

\fi

\newpage

Please note the quiz is based on the following topics:
\begin{itemize}
	\item Homomorphisms.
	\item Isomorphisms.
\end{itemize}

\begin{problem}
	Complete the following two definitions. For each, let $G$ and $H$ be two groups, and let $\phi \colon G \rightarrow H$ be a function.
\end{problem}

\begin{subproblem}
	$\phi$ is called a(n) \underline{\hspace{6em}} if $\phi(xy) = \phi(x)\phi(y)$ \underline{\hspace{2em}} $x, y \in G$.
\end{subproblem}

\begin{solution}
	$\phi$ is called a \underline{homomorphism} if $\phi(xy) = \phi(x)\phi(y)$ \underline{for all} $x, y \in G$.
\end{solution}

\begin{subproblem}
	$\phi$ is called a(n) \underline{\hspace{6em}} if it is both \underline{\hspace{2em}} and \underline{\hspace{6em}}.
\end{subproblem}

\begin{solution}
	$\phi$ is called an \underline{isomorphism} if it is both \underline{a bijection} and \underline{a homomorphism}.
\end{solution}

\begin{problem}
	Define a homomorphism $\phi \colon \mathbb{Z}_{6} \rightarrow \mathbb{Z}_{6}$ by $\phi(x) = 3x$. Which elements of $\mathbb{Z}_{6}$ belong to $\ker(\phi)$?
\end{problem}

\begin{solution}
	The set of all elements in $\ker(\phi)$ is $\set{0, 2, 4}$. \\

	By the definition of the kernel of a homomorphism, we know that an element $x \in \mathbb{Z}_{6}$ is in $\ker(\phi)$ if $\phi(x)$ equals the identity element of $\mathbb{Z}_{6}$. Taking $\mathbb{Z}_{6}$ to be an additive group, we have that $x \in \ker(\phi)$ if $\phi(x) = 0$. Using the given definition of $\phi$ and a lil' elbow grease, we can see that because $\mathbb{Z}_{6} = \set{0, 1, 2, 3, 4, 5}$, we have that:

	\vspace{2em}

	\begin{tabularx}{\linewidth}{Y Y}
		$\boxed{\phi(0) = (3)(0) \equiv 0 \mod{6}}$ & $\phi(3) = (3)(3) \equiv 3 \mod{6}$ \\\\
		$\phi(1) = (3)(1) \equiv 3 \mod{6}$ & $\boxed{\phi(3) = (3)(4) \equiv 0 \mod{6}}$ \\\\
		$\boxed{\phi(2) = (3)(2) \equiv 0 \mod{6}}$ & $\phi(3) = (3)(5) \equiv 3 \mod{6}$ \\\\
	\end{tabularx}

	\vspace{2em}

	And thus, $\ker(\phi) = \set{0, 2, 4}$.
\end{solution}

\begin{problem}
	Define a homomorphism $\psi \colon \mathbb{Z}^{\times}_{10} \rightarrow \mathbb{Z}^{\times}_{10}$ by $\psi(x) = x^{2}$. Which elements of $\mathbb{Z}^{\times}_{10}$ belong to $\ker(\psi)$?
\end{problem}

\begin{solution}
	$\ker(\psi) = \set{1, 9}$. \\

	Using similar logic to before, and leveraging that $\mathbb{Z}^{\times}_{10} = \set{1, 3, 7, 9}$ and $\psi(x) = x^{2}$, we can see that since:

	\vspace{2em}

	\begin{tabularx}{\linewidth}{Y Y}
		$\boxed{\psi(1) = 1 \equiv 1 \mod{10}}$ & $\psi(7) = 49 \equiv 9 \mod{10}$ \\\\
		$\psi(2) = 4 \equiv 4 \mod{10}$ & $\boxed{\psi(9) = 81 \equiv 1 \mod{10}}$ \\\\
	\end{tabularx}

	\vspace{2em}

	and the identity element of $\mathbb{Z}^{\times}_{10}$ is 1, we can conclude that $\ker(\psi) = \set{1, 9}$.
\end{solution}

\begin{problem}
	Let $\phi \colon G \rightarrow H$ be a group homomorphism. Which of the following statements are true?
\end{problem}

\begin{subproblem}
	$\ker(\phi)$ is a normal subgroup of $H$.
\end{subproblem}

\begin{solution}
	This statement is false, and likely derives from a fundamental misunderstanding of what the kernel even is. \\

	Recall that $\ker(\phi)$ is defined by $\ker(\phi) = \set{g \in G \colon \phi(g) = 1}$. Notice that this means every element in $\ker(\phi)$ is an element of $G$, not $H$! So clearly, it isn't \underline{always} the case $ $
\end{solution}


































\ifdefined\iscombined
	\newpage
\else
	\end{document}
\fi
